\chapter{Lock Design --- Spinlocks, MCS Locks, Futexes, and Contention}

``Use a mutex'' is the right default and a non-answer at the systems level, because a mutex hides a design space of radically different cost and fairness. A spinlock burns CPU but never sleeps; a futex sleeps but pays a syscall; a ticket lock is fair but bounces cache lines; an MCS lock scales but needs a per-thread node. This chapter builds locks from the atomic up, exposes the coherence cost that makes naive locks collapse, and gives the cost model for choosing.

\section*{Chapter Roadmap}
\begin{itemize}
\item 95.1 What a Lock Must Do, and the Cost of Contention
\item 95.2 The Naive Spinlock and Why It Collapses
\item 95.3 Test-and-Test-and-Set with Back-off
\item 95.4 Ticket Locks: Fairness
\item 95.5 MCS Locks: Scalable Queueing
\item 95.6 Futexes: Sleeping Without a Syscall (Usually)
\item 95.7 Adaptive and Reader-Writer Locks
\item 95.8 Choosing a Lock --- and Choosing Not To
\end{itemize}

\section{What a Lock Must Do, and the Cost of Contention}
A lock provides mutual exclusion. Design axes: spin or sleep, fair (FIFO) or arbitrary, and behaviour as contenders grow. The dominant cost under contention is cache coherence: the lock variable is a shared line, and every contender that writes it forces a coherence round-trip.

\textbf{Why this matters / cost model.} An uncontended lock is one atomic RMW. A contended lock is dominated by the line bouncing: N spinners writing the line generate O(N) traffic per release. Lock design is largely minimising writes to the shared line under contention.

\section{The Naive Spinlock and Why It Collapses}
\begin{lstlisting}[language=C++, caption={A test-and-set spinlock; every spin writes the shared line. Min standard: C++11.}]
class SpinLock {
    std::atomic<bool> locked_{false};
public:
    void lock() { while (locked_.exchange(true, std::memory_order_acquire)) {} }
    void unlock() { locked_.store(false, std::memory_order_release); }
};
\end{lstlisting}

\textbf{Why this matters.} Fine uncontended, but \texttt{exchange} is a write, so every spinner invalidates every copy each iteration; the line ping-pongs, the holder is slowed releasing, and throughput drops as cores are added. Burns 100\% CPU and is unfair. Do not ship a naive spinlock.

\section{Test-and-Test-and-Set with Back-off}
\begin{lstlisting}[language=C++, caption={TTAS: spin on a read, write only when free; PAUSE is x86-specific. Min standard: C++11.}]
class TTASLock {
    std::atomic<bool> locked_{false};
public:
    void lock() {
        for (;;) {
            while (locked_.load(std::memory_order_relaxed)) cpu_relax();   // read-spin: no traffic
            if (!locked_.exchange(true, std::memory_order_acquire)) return;
        }
    }
    void unlock() { locked_.store(false, std::memory_order_release); }
    static void cpu_relax() {
    #if defined(__x86_64__)
        __builtin_ia32_pause();
    #else
        std::this_thread::yield();
    #endif
    }
};
\end{lstlisting}

\textbf{Why this matters / cost model.} The read-spin keeps the line Shared --- no coherence traffic while spinning; traffic occurs only at release and the brief write race. \texttt{PAUSE} reduces power and frees pipeline/SMT resources. Better than the naive lock, but still unfair and bounces the line once per release; true scaling needs queueing.

\section{Ticket Locks: Fairness}
\begin{lstlisting}[language=C++, caption={A ticket lock: FIFO-fair, but all waiters spin on serving_. Min standard: C++11.}]
class TicketLock {
    std::atomic<unsigned> next_{0}, serving_{0};
public:
    void lock() {
        unsigned my = next_.fetch_add(1, std::memory_order_relaxed);
        while (serving_.load(std::memory_order_acquire) != my) TTASLock::cpu_relax();
    }
    void unlock() { serving_.fetch_add(1, std::memory_order_release); }
};
\end{lstlisting}

\textbf{Why this matters / cost model.} Fairness prevents starvation (a tail-latency disaster). But all waiters spin on one \texttt{serving\_}, so every release invalidates every copy --- O(N) traffic per release. Fair but not scalable; for both, each waiter must spin on a different line (MCS).

\section{MCS Locks: Scalable Queueing}
\begin{lstlisting}[language=C++, caption={MCS lock: each waiter spins on its own node; O(1) traffic per handoff. Min standard: C++11.}]
struct McsNode { std::atomic<McsNode*> next{nullptr}; std::atomic<bool> locked{true}; };
class McsLock {
    std::atomic<McsNode*> tail_{nullptr};
public:
    void lock(McsNode* self) {
        self->next.store(nullptr, std::memory_order_relaxed);
        self->locked.store(true, std::memory_order_relaxed);
        McsNode* pred = tail_.exchange(self, std::memory_order_acq_rel);
        if (pred) {
            pred->next.store(self, std::memory_order_release);
            while (self->locked.load(std::memory_order_acquire)) ;
        }
    }
    void unlock(McsNode* self) {
        McsNode* succ = self->next.load(std::memory_order_acquire);
        if (!succ) {
            McsNode* expected = self;
            if (tail_.compare_exchange_strong(expected, nullptr,
                    std::memory_order_release, std::memory_order_relaxed)) return;
            while (!(succ = self->next.load(std::memory_order_acquire))) ;
        }
        succ->locked.store(false, std::memory_order_release);
    }
};
\end{lstlisting}

\textbf{Why this matters / cost model.} Each waiter spins on a private node, so a release writes exactly one other line --- O(1) traffic per handoff regardless of contender count, scaling to dozens of cores. FIFO-fair. Cost: each acquirer supplies a node, the protocol is complex, and uncontended acquisition is slightly heavier. Linux queued spinlocks use MCS-family designs.

\section{Futexes: Sleeping Without a Syscall (Usually)}
\begin{lstlisting}[language=C++, caption={Futex fast/slow path; std::mutex is built on this. Linux-specific. Min standard: C++11.}]
// lock():   if (CAS 0->1) return;              // fast path, no syscall
//           else futex_wait(&val, 1);          // slow path, sleep in kernel
// unlock(): store 0; if (waiters) futex_wake(&val, 1);
\end{lstlisting}

\textbf{Why this matters / cost model.} The uncontended case never enters the kernel (one CAS); only contention triggers \texttt{futex\_wait}/\texttt{futex\_wake} (syscall + context switch). \texttt{std::mutex} is built on this. Sleeping is good for long or rare-contention sections and oversubscription but catastrophic for sub-microsecond sections; spin there instead.

\section{Adaptive and Reader-Writer Locks}
Adaptive mutexes spin briefly then futex-sleep --- the sensible default that good \texttt{std::mutex} implementations use. Reader-writer locks (\texttt{std::shared\_mutex}) allow many readers or one writer, but the shared reader count is itself a contended atomic; under short frequent reads the rwlock can be slower than a plain mutex --- exactly where RCU wins.

\textbf{Why this matters.} A modern \texttt{std::mutex} already spins-then-sleeps optimally, so it is good default advice. Use an rwlock only for long read sections; for read-mostly short sections prefer RCU/seqlocks.

\section{Choosing a Lock --- and Choosing Not To}
\begin{itemize}
\item Naive spinlock --- never ship.
\item TTAS+back-off --- short sections, low contention.
\item Ticket lock --- short sections needing fairness; does not scale.
\item MCS lock --- short sections, high core count; scales and is fair.
\item Futex/\texttt{std::mutex} --- general; long or rare-contention sections.
\item Adaptive mutex --- the sensible default.
\item \texttt{shared\_mutex} --- long read sections only.
\end{itemize}

\textbf{The discipline.} Choose by critical-section length and contention: sub-microsecond, lightly contended $\to$ spin (TTAS/MCS); longer or oversubscribed $\to$ sleep (futex); default to adaptive \texttt{std::mutex}. The deepest lesson: the best lock is often no lock --- partition per thread (Chapter 96), use an SPSC ring (Chapter 77), or make reads free with RCU (Chapter 94).
